\chapter{Implementation Instructions}
\label{ch:implementation-instructions}

This chapter provides practical guidance for implementing and maintaining
the HITL Requirements Traceability framework, including tooling setup,
validation procedures, and maintenance workflows.

% =============================================================================
\section{Document Compilation}
\label{sec:document-compilation}

\subsection{Prerequisites}

\begin{itemize}[noitemsep]
  \item \textbf{TeX Distribution:} TeX Live 2024+ or MiKTeX 24+
  \item \textbf{Engine:} XeLaTeX (required for fontspec)
  \item \textbf{Packages:} Standard LaTeX packages (included in full distributions)
\end{itemize}

\subsection{Compilation Command}

\begin{lstlisting}[language=bash,basicstyle=\ttfamily\small]
cd docs/latex/specs/tb-hitl
xelatex -interaction=nonstopmode tb-hitl-traceability-spec.tex
xelatex -interaction=nonstopmode tb-hitl-traceability-spec.tex
\end{lstlisting}

\paragraph{Note:} Two passes are required for cross-references and
table of contents generation.

\subsection{Build Integration}

Add to the project Makefile:

\begin{lstlisting}[language=make,basicstyle=\ttfamily\small]
tb-hitl-spec:
	cd docs/latex/specs/tb-hitl && \
	xelatex -interaction=nonstopmode tb-hitl-traceability-spec.tex && \
	xelatex -interaction=nonstopmode tb-hitl-traceability-spec.tex

tb-hitl-pdf: tb-hitl-spec
	cp docs/latex/specs/tb-hitl/tb-hitl-traceability-spec.pdf \
	   docs/pdf/specs/tb-hitl-traceability-spec.pdf
\end{lstlisting}

% =============================================================================
\section{Traceability Link Validation}
\label{sec:link-validation}

\subsection{Manual Validation}

For each traceability link, verify:

\begin{enumerate}
  \item \textbf{Source exists:} The referenced source document is present
        at the specified path
  \item \textbf{Target exists:} The link target (requirement, spec, schema)
        is defined in its respective chapter
  \item \textbf{Bidirectional:} The reverse link is also documented
  \item \textbf{Consistent:} Link metadata matches between forward and
        backward traces
\end{enumerate}

\subsection{Automated Validation Script}

Create a Python validation script at \texttt{scripts/validate-hitl-links.py}:

\begin{lstlisting}[language=python,basicstyle=\ttfamily\small]
#!/usr/bin/env python3
"""Validate HITL traceability links."""

import yaml
import json
from pathlib import Path

def load_manifest():
    """Load the HITL traceability manifest."""
    manifest_path = Path("docs/tb/machine-readable/hitl-traceability-manifest.yaml")
    with open(manifest_path) as f:
        return yaml.safe_load(f)

def validate_source_exists(manifest):
    """Check all source documents exist."""
    errors = []
    for source in manifest.get("source_documents", []):
        path = Path(source["filepath"])
        if not path.exists():
            errors.append(f"Source not found: {source['source_ref']}")
    return errors

def validate_forward_links(manifest):
    """Validate forward traceability links."""
    errors = []
    requirements = {r["requirement_id"] for r in manifest.get("requirements", [])}
    
    for link in manifest.get("forward_traceability", []):
        if link["requirement_id"] not in requirements:
            errors.append(f"Unknown requirement: {link['requirement_id']}")
    return errors

if __name__ == "__main__":
    manifest = load_manifest()
    
    errors = []
    errors.extend(validate_source_exists(manifest))
    errors.extend(validate_forward_links(manifest))
    
    if errors:
        print("Validation FAILED:")
        for e in errors:
            print(f"  - {e}")
        exit(1)
    else:
        print("Validation PASSED")
        exit(0)
\end{lstlisting}

% =============================================================================
\section{Adding New Source Documents}
\label{sec:adding-sources}

When a new source document is added to the TB corpus:

\begin{enumerate}
  \item \textbf{Catalogue the document:}
        \begin{itemize}[noitemsep]
          \item Assign a source reference ID (e.g., \srcref{TB-NEW-001})
          \item Add entry to \cref{ch:source-documents}
          \item Update \texttt{hitl-traceability-manifest.yaml}
        \end{itemize}
        
  \item \textbf{Extract to machine-readable format:}
        \begin{itemize}[noitemsep]
          \item Run appropriate extraction script
          \item Store output in \texttt{docs/tb/machine-readable/extracted/}
          \item Add YAML frontmatter with source reference
        \end{itemize}
        
  \item \textbf{Derive requirements:}
        \begin{itemize}[noitemsep]
          \item Review extracted content for requirement statements
          \item Assign requirement IDs following naming convention
          \item Document derivation evidence in \cref{ch:business-requirements}
        \end{itemize}
        
  \item \textbf{Update traceability matrices:}
        \begin{itemize}[noitemsep]
          \item Add forward links in \cref{ch:bidirectional-traceability}
          \item Add backward links from affected specifications
          \item Run validation script to verify links
        \end{itemize}
        
  \item \textbf{Update gap analysis:}
        \begin{itemize}[noitemsep]
          \item Check for new gaps introduced
          \item Update orphan analysis tables
          \item Document any new gaps in \cref{ch:gap-analysis}
        \end{itemize}
\end{enumerate}

% =============================================================================
\section{Adding New Requirements}
\label{sec:adding-requirements}

\subsection{Requirement ID Convention}

\begin{table}[h]
\centering
\caption{Requirement ID Naming Convention}
\begin{tabularx}{\textwidth}{@{}L{3.5cm}L{4cm}L{6cm}@{}}
\toprule
\textbf{Pattern} & \textbf{Category} & \textbf{Example} \\
\midrule
TB-REQ-ES\textit{nn} & Executive Summary & TB-REQ-ES01 \\
TB-REQ-SC\textit{nn} & Scope & TB-REQ-SC01 \\
TB-REQ-PS\textit{nn} & Process Steps & TB-REQ-PS01 \\
TB-REQ-AI\textit{nn} & AI Augmentation & TB-REQ-AI01 \\
BSS-REQ-\textit{nn} & BSS Risk Assessment & BSS-REQ-01 \\
TB-REQ-FI\textit{nn} & Financial (future) & TB-REQ-FI01 \\
\bottomrule
\end{tabularx}
\end{table}

\subsection{Requirement Documentation Template}

\begin{lstlisting}[language=tex,basicstyle=\ttfamily\small]
\subsection{\reqref{TB-REQ-XX01}: Requirement Title}
\label{sec:req-xx01}

\begin{table}[h]
\centering
\caption{\reqref{TB-REQ-XX01} Derivation Evidence}
\begin{tabularx}{\textwidth}{@{}L{3.5cm}L{10.5cm}@{}}
\toprule
\textbf{Attribute} & \textbf{Value} \\
\midrule
Requirement ID & TB-REQ-XX01 \\
Requirement Text & The system shall... \\
Source Document & \srcref{TB-XXX-001}: Document name \\
Source Location & Section X, Paragraph Y \\
Derivation Method & Direct extraction / Synthesis \\
Confidence & High / Medium / Low \\
Status & Draft / Approved \\
\bottomrule
\end{tabularx}
\end{table}

\paragraph{Source Quote:}
\begin{quote}
\textit{``Quoted text from source document...''}
\end{quote}
\end{lstlisting}

% =============================================================================
\section{Schema Evolution}
\label{sec:schema-evolution}

When schemas are modified:

\begin{enumerate}
  \item \textbf{Document the change:}
        \begin{itemize}[noitemsep]
          \item Update schema file with new/modified fields
          \item Increment schema version
          \item Add changelog entry
        \end{itemize}
        
  \item \textbf{Update traceability:}
        \begin{itemize}[noitemsep]
          \item Add new fields to \cref{ch:schema-mapping}
          \item Link fields to specifications and requirements
          \item Verify no orphan fields remain
        \end{itemize}
        
  \item \textbf{Assess impact:}
        \begin{itemize}[noitemsep]
          \item Identify affected downstream consumers
          \item Document migration requirements
          \item Update validation scripts
        \end{itemize}
\end{enumerate}

% =============================================================================
\section{HITL Review Execution}
\label{sec:review-execution}

\subsection{Pre-Review Checklist}

Before initiating HITL review:

\begin{itemize}
  \item[$\square$] Document compiled successfully with no LaTeX errors
  \item[$\square$] All chapter files present and included
  \item[$\square$] Validation script passes
  \item[$\square$] Machine-readable manifest is current
  \item[$\square$] Reviewers identified and available
\end{itemize}

\subsection{Review Meeting Agenda}

\begin{enumerate}
  \item \textbf{Introduction (5 min):} Review objectives and scope
  \item \textbf{Source Review (15 min):} Checkpoint H1 --- verify source completeness
  \item \textbf{Extraction Review (15 min):} Checkpoint H2 --- sample extraction accuracy
  \item \textbf{Requirements Review (20 min):} Checkpoint H3 --- verify derivations
  \item \textbf{Traceability Review (15 min):} Checkpoint H4 --- validate links
  \item \textbf{Schema Review (15 min):} Checkpoint H5 --- verify alignment
  \item \textbf{Gap Review (10 min):} Checkpoint H6 --- accept gap plans
  \item \textbf{Sign-Off (5 min):} Complete sign-off register
\end{enumerate}

\subsection{Post-Review Actions}

\begin{enumerate}
  \item Document review findings in sign-off register
  \item Create issues for any identified defects
  \item Update gap register if new gaps identified
  \item Publish approved version to document repository
  \item Notify stakeholders of approval status
\end{enumerate}

% =============================================================================
\section{Continuous Integration}
\label{sec:ci-integration}

\subsection{CI Pipeline Configuration}

Add to \texttt{.github/workflows/docs.yml}:

\begin{lstlisting}[language=yamlx,basicstyle=\ttfamily\small]
name: HITL Documentation

on:
  push:
    paths:
      - 'docs/latex/specs/tb-hitl/**'
      - 'docs/tb/machine-readable/**'

jobs:
  build-and-validate:
    runs-on: ubuntu-latest
    steps:
      - uses: actions/checkout@v4
      
      - name: Install TeX Live
        run: |
          sudo apt-get update
          sudo apt-get install -y texlive-xetex texlive-latex-extra
      
      - name: Validate traceability links
        run: python scripts/validate-hitl-links.py
      
      - name: Compile specification
        run: |
          cd docs/latex/specs/tb-hitl
          xelatex -interaction=nonstopmode tb-hitl-traceability-spec.tex
          xelatex -interaction=nonstopmode tb-hitl-traceability-spec.tex
      
      - name: Upload PDF artifact
        uses: actions/upload-artifact@v4
        with:
          name: tb-hitl-spec
          path: docs/latex/specs/tb-hitl/tb-hitl-traceability-spec.pdf
\end{lstlisting}

% =============================================================================
\section{Troubleshooting}
\label{sec:troubleshooting}

\begin{table}[h]
\centering
\caption{Common Issues and Solutions}
\begin{tabularx}{\textwidth}{@{}L{5cm}L{8.5cm}@{}}
\toprule
\textbf{Issue} & \textbf{Solution} \\
\midrule
LaTeX compilation error & Check for missing packages; verify fontspec compatibility \\
Undefined cross-reference & Run xelatex twice; check label spelling \\
Missing source document & Verify file path in manifest; check for renamed files \\
Broken traceability link & Run validation script; check ID spelling \\
Table overflow & Reduce column widths; use longtable for long tables \\
TikZ positioning error & Verify tikzlibrary imports; check node distance \\
\bottomrule
\end{tabularx}
\end{table}